\documentclass[11pt]{article}
\usepackage[margin=1in]{geometry}
\usepackage[T1]{fontenc}
\usepackage[utf8]{inputenc}
\usepackage{lmodern}
\usepackage{microtype}
\usepackage{amsmath,amssymb}
\usepackage{booktabs}
\usepackage{graphicx}
\usepackage{xcolor}
\usepackage[colorlinks=true,linkcolor=blue!50!black,citecolor=blue!50!black,urlcolor=blue!50!black]{hyperref}

\title{Structure Demotion: Reclaiming the KV-Cache Budget from Repeated Separators under Extreme Compression}
\author{Satyajeeth Suresh Kannan\\
Innovation Academy, Alpharetta, Georgia, USA\\
\texttt{jeetusk08@gmail.com}
}
\date{}

\emergencystretch=3em
\begin{document}
\maketitle

\begin{abstract}
Query-agnostic KV-cache eviction methods such as KVzip score every cached token by how much it helps the model reconstruct its own
context. At extreme compression this criterion misallocates the budget: repeated structural tokens---newlines, list markers,
separators---are easy to reconstruct from one another and keep a large share of the cache, while repeated content words, exactly
the evidence that aggregation tasks need, are discounted. We propose \emph{structure demotion}, a training-free rule that multiplies
the scores of structural tokens (tokens containing no letter or digit) that occur at least four times in the context by $0.25$
before eviction, while exempting ``joiners'' inside identifiers such as the dashes of a UUID. Applied to RestoreKV+, the strongest
entry on the kvpress leaderboard, at compression ratio $0.9375$ (one sixteenth of the cache kept) on RULER-4096 with Qwen3-8B,
structure demotion raises the mean score from 86.38 to 89.64 ($+3.25$ points; paired 95\% confidence interval $[+2.73, +3.77]$),
recovering 36\% of the gap to the uncompressed model. Seven of the thirteen tasks improve significantly and none degrades
significantly; common-words extraction gains 12.1 points. We also report an instructive failure: demoting digits, or the dashes
inside UUIDs, destroys needle retrieval. The rule is a default-off option of about thirty lines in the kvpress library.
\end{abstract}

\section{Introduction}
The key--value (KV) cache of a transformer grows linearly with context length and dominates inference memory for long inputs.
KV-cache eviction keeps only a subset of the cached keys and values. Query-agnostic methods compress the cache once, before any
question is known, so that the compressed cache can be reused across queries. KVzip~\cite{kim2025kvzip} is the reference method
in this setting: it asks the model to reconstruct its context and keeps the KV pairs that receive the most attention during that
reconstruction. Its successors refine the scoring~\cite{jegou2026kvzap,kim2026fastkvzip} or add learned restore tokens that
recover what eviction loses~\cite{baek2026restorekv}. On the public kvpress leaderboard~\cite{kvpress,kvpressleaderboard}, the
strongest method for Qwen3-8B at compression ratios $0.88$ and $0.9375$ is RestoreKV+~\cite{baek2026restorekv}.

At compression ratio $0.9375$, where only one of every sixteen cached positions survives, RestoreKV+ still loses 8.9 of the 95.3
points that Qwen3-8B~\cite{yang2025qwen3} scores on RULER-4096~\cite{hsieh2024ruler} with a full cache. The loss is
concentrated: common-words and frequent-words extraction and the multi-needle tasks account for most of it. We traced it to what
the reconstruction criterion rewards. A token that the model can easily reconstruct from its neighbours is not necessarily
unimportant, and a token the model keeps because other copies predict it is not necessarily useful. In long structured contexts
the cheapest tokens to keep ``in the loop'' of reconstruction are the repeated separators: newlines, list numbers, dashes and
quotes. They crowd out the content words the question needs, and repeated content words---which reconstruction discounts because
later copies are predictable---are exactly the answer in aggregation tasks.

We make three contributions.
\begin{itemize}
  \item A diagnosis of where the budget goes under reconstruction scoring at extreme compression (Section~\ref{sec:diagnosis}).
  \item \emph{Structure demotion}, a one-line change to the eviction scores with a single exemption rule, implemented as a
        default-off option of \texttt{KVzipPress} and \texttt{RestoreKVPress} in kvpress (Section~\ref{sec:method}).
  \item A full-benchmark evaluation with paired confidence intervals against the published leaderboard predictions, plus the
        ablations and the failure case that shaped the rule (Section~\ref{sec:experiments}).
\end{itemize}

\section{Background}
\label{sec:background}
\paragraph{Compression ratio.} Following kvpress, the compression ratio is the fraction of cached context positions that are
removed; at $0.9375$ the model keeps $6.25\%$ of the context KV pairs, averaged over layers and heads.

\paragraph{KVzip.} After prefilling the context, KVzip~\cite{kim2025kvzip} runs the model on a prompt that asks it to repeat the
context and records, for every cached KV pair, the maximum attention it receives from the reconstruction queries. Pairs are then
evicted by a global bottom-$k$ over all layers and heads, so the budget is non-uniform across layers. The KVzip+ normalization
introduced with KVzap~\cite{jegou2026kvzap} rescales these attention weights by the norm of the value vector after the output
projection and by the inverse norm of the query-side hidden state.

\paragraph{RestoreKV.} RestoreKV~\cite{baek2026restorekv} appends $n=8$ learned restore tokens that attend to the full cache in
one LoRA-adapted pass before eviction; the resulting restore cache occupies part of the same budget, and KVzip-style pruning fills
the rest. RestoreKV+ is RestoreKV with KVzip+ normalization. The restore embeddings and LoRA weights are trained by
self-distillation from the full-cache model and are published for Qwen3-8B.

\paragraph{RULER-4096.} RULER~\cite{hsieh2024ruler} contains thirteen synthetic long-context tasks: eight needle-in-a-haystack
variants (single, multi-key, multi-value, multi-query), variable tracking (\texttt{vt}), common-words and frequent-words extraction
(\texttt{cwe}, \texttt{fwe}) and two question-answering tasks (\texttt{qa\_1}, \texttt{qa\_2}). The kvpress benchmark uses 500
examples per task at a 4{,}096-token context and scores each task by string matching (the fraction of reference answers found in
the output; any reference for QA). The leaderboard reports the unweighted mean over tasks.

\section{Where the budget goes}
\label{sec:diagnosis}
We instrumented the eviction step of RestoreKV+ on a common-words extraction example (5{,}127 tokens; the model must list the ten
most frequent words in a numbered list of words). Table~\ref{tab:budget} shows how the $6.04\%$ of the cache that survives is
spent. List-index numbers make up 60\% of the context and take 38\% of the budget; whitespace and newline tokens make up 16\% of
the context and are kept at almost the same rate as words (8.7\% versus 10.0\%), taking 23\% of the budget. Only 38\% of the budget
goes to words.

The reconstruction criterion also discounts repetition. Averaged over the words of this example, the first three occurrences of a
word score about $0.045$ and are kept at a rate of 12.7\%; from the fourth occurrence on, the score halves (about $0.022$--$0.026$)
and the kept rate drops to 5.3--6.5\%. A fourth copy of a word is predictable from the first three, so reconstruction needs it less.
On a counting task this is backwards: each of the ten gold words occurs 30 times, and their mean score ($0.026$) is below the mean
score of all word tokens ($0.037$). Only 6.9\% of the gold-word occurrences survive.

\begin{table}[t]
\centering
\small
\begin{tabular}{lrrrrr}
\toprule
 & & \multicolumn{2}{c}{RestoreKV+} & \multicolumn{2}{c}{+ structure demotion} \\
\cmidrule(lr){3-4}\cmidrule(lr){5-6}
token class & share of context & kept & share of budget & kept & share of budget \\
\midrule
words & 23.0\% & 10.0\% & 38\% & 13.0\% & 50\% \\
numbers (list indices) & 60.4\% & 3.8\% & 38\% & 4.3\% & 43\% \\
whitespace / newlines & 16.2\% & 8.7\% & 23\% & 2.4\% & 6\% \\
other punctuation & 0.4\% & 10.7\% & 1\% & 13.6\% & 1\% \\
\midrule
gold-word occurrences kept & & \multicolumn{2}{c}{6.9\%} & \multicolumn{2}{c}{9.7\%} \\
score on this example & & \multicolumn{2}{c}{90} & \multicolumn{2}{c}{100} \\
\bottomrule
\end{tabular}
\caption{Budget allocation on one common-words extraction example (Qwen3-8B, compression ratio $0.9375$, $6.04\%$ of the context
kept in both cases). Structure demotion moves budget from newlines to words and raises the retention of the answer words by 41\%.}
\label{tab:budget}
\end{table}

\section{Structure demotion}
\label{sec:method}
Let $x_1,\dots,x_n$ be the token ids of the context and $s_{\ell,h,i}$ the eviction score of position $i$ in layer $\ell$ and KV head
$h$. We call a token id \emph{structural} if its decoded string, with surrounding whitespace removed, is not a word, word piece or
number, that is, it does not fully match the pattern \verb+[^\W_](?:[^\W_]|['\-])*+ (a letter or digit followed by letters,
digits, apostrophes or hyphens). Newlines, list punctuation, quotes, brackets and lone dashes are structural; words, word pieces,
digits and numbers are not. Position $i$ is \emph{demoted} when
\begin{enumerate}
  \item $x_i$ is structural;
  \item the id $x_i$ occurs at least $m$ times in the context (we use $m=4$, the occurrence at which the reconstruction score halves
        in Section~\ref{sec:diagnosis}); and
  \item $x_i$ is not a \emph{joiner}: a structural token without whitespace whose previous token ends with a letter or digit and whose
        next token starts with one, such as the dashes of a UUID, the point of \texttt{3.14}, the colon of \texttt{10:30} or the comma
        of \texttt{1,000}.
\end{enumerate}
Before the global bottom-$k$ eviction we set $s_{\ell,h,i} \leftarrow \alpha\, s_{\ell,h,i}$ for every demoted position, in every
layer and head, with $\alpha = 0.25$. Positions before the end of the chat prefix and the attention-sink positions are never
demoted, and RestoreKV's restore tokens lie outside the context and are untouched. The number of evicted pairs is unchanged, so the
achieved compression ratio is identical ($6.04\%$ kept in the example of Table~\ref{tab:budget} with and without demotion).

The rule needs only the token ids and the tokenizer that KVzip already loads. It is a single pass over the context on the CPU, far
cheaper than KVzip's reconstruction passes. With $\alpha=0$ the option is disabled and every existing press is bit-identical.

\paragraph{Why joiners are exempt.} The first full-benchmark run used conditions (1) and (2) only. It scored 64.6 on
\texttt{niah\_multikey\_3}, against 97.0 for the baseline. In that task the keys and values are UUIDs, which the Qwen tokenizer
splits into letter pieces, digits and standalone dash tokens; the dash repeats hundreds of times per context, so it was demoted and
the UUIDs were broken up by eviction. Exempting joiners reduces the number of demoted positions in a typical
\texttt{niah\_multikey\_3} context from 566 to 172 and restores the task to 96.8, while leaving the demoted set unchanged on the
extraction and variable-tracking tasks. Separators that carry whitespace on at least one side (list markers, newlines, quotes) are
still demoted.

\section{Experiments}
\label{sec:experiments}
\paragraph{Setup.} We evaluate Qwen3-8B in bfloat16 on the kvpress RULER-4096 benchmark (13 tasks $\times$ 500 examples) at
compression ratio $0.9375$, using the published RestoreKV+ adapters for Qwen3-8B and the kvpress evaluation script unchanged except
for the new option (\texttt{RestoreKVPress(kvzip\_\allowbreak plus\_\allowbreak normalization=True,\allowbreak{} structure\_\allowbreak demotion=0.25)}). The run took about 35
hours on an Apple M5 Pro laptop (PyTorch MPS backend, SDPA attention). Leaderboard entries were produced on CUDA with
FlashAttention; on the first 100 common-words examples our local run of the unmodified baseline scored 77.2, in line with the
published 78.88 over all 500.

\paragraph{Paired comparison.} Because the leaderboard publishes per-example predictions, we compare against RestoreKV+ on exactly
the same 6{,}500 examples. We recompute per-example scores with the leaderboard's metric (the recomputation reproduces the
published task scores exactly for twelve tasks and within $0.2$ points for \texttt{qa\_2}) and bootstrap the per-example differences
within each task ($10{,}000$ resamples) to obtain 95\% confidence intervals for each task and for the 13-task mean.

\paragraph{Development protocol.} All design choices---the structural test, $m$, $\alpha$ and the joiner exemption---were made on
paired 100-example screens (the first 100 examples of a task) and, for the joiner exemption, after the failure of the first full
run described in Section~\ref{sec:method}. The final configuration was then run once on the full benchmark.

\begin{table}[t]
\centering
\small
\setlength{\tabcolsep}{4pt}
\begin{tabular}{lrrrcr}
\toprule
task & RestoreKV+ & + demotion & $\Delta$ & 95\% CI & better / worse \\
\midrule
cwe & 78.88 & \textbf{90.96} & $+12.08$ & $[+11.08, +13.14]$ & 362 / 12 \\
fwe & 83.67 & \textbf{90.00} & $+6.33$ & $[+5.07, +7.67]$ & 104 / 9 \\
niah\_multikey\_1 & 85.80 & \textbf{92.20} & $+6.40$ & $[+3.40, +9.40]$ & 45 / 13 \\
niah\_multikey\_2 & \textbf{99.60} & 99.20 & $-0.40$ & $[-1.20, +0.40]$ & 1 / 3 \\
niah\_multikey\_3 & \textbf{97.00} & 96.80 & $-0.20$ & $[-1.40, +1.00]$ & 4 / 5 \\
niah\_multiquery & 96.80 & \textbf{98.20} & $+1.40$ & $[+0.55, +2.25]$ & 43 / 14 \\
niah\_multivalue & 82.75 & \textbf{88.15} & $+5.40$ & $[+3.70, +7.15]$ & 148 / 55 \\
niah\_single\_1 & 100.00 & 100.00 & $0.00$ & $[0.00, 0.00]$ & 0 / 0 \\
niah\_single\_2 & 90.80 & \textbf{96.00} & $+5.20$ & $[+2.80, +7.60]$ & 32 / 6 \\
niah\_single\_3 & 93.00 & \textbf{95.80} & $+2.80$ & $[+0.40, +5.20]$ & 27 / 13 \\
qa\_1 & 65.40 & \textbf{67.60} & $+2.20$ & $[-0.60, +5.00]$ & 32 / 21 \\
qa\_2 & 49.20 & \textbf{50.40} & $+1.20$ & $[-2.00, +4.00]$ & 33 / 28 \\
vt & 100.00 & 100.00 & $0.00$ & $[0.00, 0.00]$ & 0 / 0 \\
\midrule
mean & 86.38 & \textbf{89.64} & $+3.25$ & $[+2.73, +3.77]$ & \\
\bottomrule
\end{tabular}
\caption{RULER-4096, Qwen3-8B, compression ratio $0.9375$. RestoreKV+ scores are the published leaderboard entry; confidence
intervals and example counts come from the paired per-example comparison (the $\Delta$ for \texttt{qa\_2} is $+1.00$ against the
recomputed baseline of 49.40). For reference, the uncompressed model scores 95.32 and RestoreKV without KVzip+ normalization 81.93
at the same ratio.}
\label{tab:main}
\end{table}

\paragraph{Main result.} Table~\ref{tab:main} shows the full comparison. Structure demotion raises the mean from 86.38 to 89.64, a
gain of $3.25$ points with a paired 95\% confidence interval of $[+2.73, +3.77]$ (no bootstrap resample reached zero). It closes
36\% of the gap between RestoreKV+ and the uncompressed model (95.32). Seven tasks improve significantly; the largest gains are on
the two extraction tasks (\texttt{cwe} $+12.1$, \texttt{fwe} $+6.3$) and on the needle tasks whose haystacks are lists or repeated
sentences (\texttt{niah\_multikey\_1}, \texttt{niah\_single\_2}, \texttt{niah\_multivalue}, each above $+5$). The two UUID tasks
move by $-0.4$ and $-0.2$, both well inside their intervals, and no task is significantly worse. The QA tasks, whose contexts are
prose with little repeated structure, change little.

\begin{table}[t]
\centering
\small
\begin{tabular}{lrrrrrrr}
\toprule
variant (100-example screens) & cwe & fwe & mk\_1 & mv & vt & qa\_1 & qa\_2 \\
\midrule
RestoreKV+ (local baseline) & 77.2 & 82.3 & 90.0 & 80.8 & 100.0 & 67.0 & 60.0 \\
demote structural tokens, $\alpha=0.25$ (ours) & 90.6 & 90.0 & 94.0 & 85.5 & 100.0 & 64.0 & 56.0 \\
demote structural tokens, $\alpha=0.5$ & 87.4 & -- & -- & -- & -- & -- & 57.0 \\
also demote repeated digits & 96.1 & 90.0 & 0.0 & 0.0 & -- & -- & -- \\
demote repeated non-alphabetic bigrams ($\geq 16$) & 90.0 & -- & 90.0 & -- & -- & -- & -- \\
boost repeated content words instead & 58.0 & -- & -- & -- & -- & -- & -- \\
\bottomrule
\end{tabular}
\caption{Development screens on the first 100 examples of each task (Qwen3-8B, ratio $0.9375$; mk\_1 = niah\_multikey\_1,
mv = niah\_multivalue; -- = not run). The joiner exemption, added later, leaves the demoted set unchanged on cwe, fwe and vt and changes at most 12 positions per context on the others. With 100
examples a single task score has a standard error of several points; the QA losses seen here did not reappear on the full
benchmark (Table~\ref{tab:main}).}
\label{tab:ablation}
\end{table}

\paragraph{Ablations.} Table~\ref{tab:ablation} summarises the development screens. Halving the demotion ($\alpha=0.5$) gives back
3.2 points on \texttt{cwe} for 1 point on \texttt{qa\_2}, so we kept $\alpha=0.25$. Demoting repeated digits as well recovers
\texttt{cwe} almost to the uncompressed level (96.1) but scores zero on the needle tasks: the Qwen tokenizer splits every number
into single digits, so a seven-digit needle is seven tokens that each ``repeat'', and losing one of them makes the answer wrong. A
rule based on repeated non-alphabetic bigrams (for example \texttt{"2"~"."}) protected needles but gained less on \texttt{cwe}.
The opposite intervention---undoing the discount on repeated content words---made things worse ($-19$ on \texttt{cwe}): the
repeated tokens it boosted were mostly the list-index digits, which moved budget in the wrong direction. The lesson is that any
repetition-based rule must be checked on needle tasks, whose answers are made of digits and identifier pieces.

\section{Related work}
\paragraph{Attention-based eviction.} H2O keeps ``heavy-hitter'' tokens with large accumulated attention~\cite{zhang2023h2o};
StreamingLLM keeps attention sinks and a recent window~\cite{xiao2024streamingllm}; SnapKV selects positions using the attention of
an observation window at the end of the prompt~\cite{li2024snapkv}; PyramidKV allocates smaller budgets to higher
layers~\cite{cai2024pyramidkv}; TOVA keeps the tokens with the highest current attention~\cite{oren2024tova}. These methods score
tokens with the queries of the prompt itself, which ties the compressed cache to one question.

\paragraph{Query-agnostic compression.} KVzip replaces the question with a context-reconstruction task~\cite{kim2025kvzip}, Fast
KVzip approximates it with a lightweight gate~\cite{kim2026fastkvzip}, KVzap approximates KVzip+ scores with a fast surrogate
model~\cite{jegou2026kvzap}, and RestoreKV adds learned restore tokens~\cite{baek2026restorekv}. Structure demotion is orthogonal
to all of them: it post-processes whatever reconstruction scores they produce.

\paragraph{Separators.} SepLLM observes that separator tokens attract disproportionate attention and trains models that compress
each segment into its separator~\cite{chen2024sepllm}. Our finding is complementary: in an off-the-shelf model scored by
reconstruction, \emph{repeated} separators are retained far beyond their value for the downstream task, and moving their share of
an extreme budget to content tokens helps.

\section{Limitations}
We evaluated one model, one context length and one compression ratio at full scale, on one benchmark whose contexts are synthetic
and highly structured; the gains are largest exactly where contexts repeat separators, and we expect little effect on plain prose.
The run used Apple's MPS backend with SDPA attention rather than the CUDA and FlashAttention setup of the other leaderboard entries;
the baseline reproduced within two points locally, and a CUDA re-run needs only the released option. The factor $\alpha$ and the
threshold $m$ were set on small screens and not tuned further. The structural test is a tokenizer-level heuristic: in code, JSON or
languages where punctuation carries meaning, demoting repeated symbols may remove information the task needs, and tokenizers that
merge punctuation with letters would see fewer structural tokens. Finally, the joiner exemption was added after observing the
failure of the first full run, so the full-benchmark number evaluates a rule that was revised once in response to that benchmark.

\section{Conclusion}
Reconstruction scoring asks which tokens the model needs in order to reproduce its context; at extreme compression the answer is
dominated by repeated scaffolding. Demoting repeated structural tokens by a constant factor, while protecting punctuation inside
identifiers, reallocates that budget to content and improves the strongest query-agnostic method on RULER-4096 by 3.25 points at
one-sixteenth of the cache. The rule is training-free, model-free and a few lines long, which makes it easy to combine with future
reconstruction-based presses.

\section*{Code and data}
The implementation is the \texttt{structure\_demotion} option of \texttt{KVzipPress}/\texttt{RestoreKVPress}, submitted to kvpress
as pull request NVIDIA/kvpress\#288 (branch \texttt{structure-demotion-restorekv} of \url{https://github.com/JeetuSK0808/kvpress}).
Predictions, metrics and configuration for the full run are in the leaderboard submission
\url{https://huggingface.co/spaces/nvidia/kvpress-leaderboard/discussions/24}.

\section*{Use of AI assistance}
The experiments, analysis and code were developed with an AI coding assistant (Claude Code, Anthropic) directed by the author, who
reviewed the results and the manuscript.

\section*{Acknowledgments}
We thank the kvpress maintainers for the benchmark, the leaderboard and its published predictions, and the RestoreKV authors for
releasing their adapters.

\begin{thebibliography}{99}

\bibitem{kim2025kvzip}
Jang-Hyun Kim, Jinuk Kim, Sangwoo Kwon, Jae W. Lee, Sangdoo Yun, and Hyun Oh Song.
\newblock KVzip: Query-Agnostic KV Cache Compression with Context Reconstruction.
\newblock arXiv preprint arXiv:2505.23416, 2025. \url{https://arxiv.org/abs/2505.23416}

\bibitem{jegou2026kvzap}
Simon Jegou and Maximilian Jeblick.
\newblock KVzap: Fast, Adaptive, and Faithful KV Cache Pruning.
\newblock arXiv preprint arXiv:2601.07891, 2026. \url{https://arxiv.org/abs/2601.07891}

\bibitem{kim2026fastkvzip}
Jang-Hyun Kim, Dongyoon Han, and Sangdoo Yun.
\newblock Fast KVzip: Efficient and Accurate LLM Inference with Gated KV Eviction.
\newblock arXiv preprint arXiv:2601.17668, 2026. \url{https://arxiv.org/abs/2601.17668}

\bibitem{baek2026restorekv}
Changwoo Baek, Seungjun Shin, and Kyeongbo Kong.
\newblock RestoreKV: Recovering Full-Cache Behavior Under Aggressive Query-Agnostic KV Cache Eviction.
\newblock arXiv preprint arXiv:2608.01247, 2026. \url{https://arxiv.org/abs/2608.01247}

\bibitem{kvpress}
NVIDIA.
\newblock kvpress: {LLM} {KV} cache compression made easy.
\newblock \url{https://github.com/NVIDIA/kvpress}, 2026. Software, accessed September 2026.

\bibitem{kvpressleaderboard}
NVIDIA.
\newblock {kvpress} leaderboard.
\newblock \url{https://huggingface.co/spaces/nvidia/kvpress-leaderboard}, 2026. Accessed September 2026.

\bibitem{yang2025qwen3}
An Yang, Anfeng Li, Baosong Yang, Beichen Zhang, Binyuan Hui, Bo Zheng, Bowen Yu, Chang Gao, et al.
\newblock Qwen3 Technical Report.
\newblock arXiv preprint arXiv:2505.09388, 2025. \url{https://arxiv.org/abs/2505.09388}

\bibitem{hsieh2024ruler}
Cheng-Ping Hsieh, Simeng Sun, Samuel Kriman, Shantanu Acharya, Dima Rekesh, Fei Jia, Yang Zhang, and Boris Ginsburg.
\newblock RULER: What's the Real Context Size of Your Long-Context Language Models?.
\newblock arXiv preprint arXiv:2404.06654, 2024. \url{https://arxiv.org/abs/2404.06654}

\bibitem{zhang2023h2o}
Zhenyu Zhang, Ying Sheng, Tianyi Zhou, Tianlong Chen, Lianmin Zheng, Ruisi Cai, Zhao Song, Yuandong Tian, et al.
\newblock {H}$_2${O}: Heavy-Hitter Oracle for Efficient Generative Inference of Large Language Models.
\newblock arXiv preprint arXiv:2306.14048, 2023. \url{https://arxiv.org/abs/2306.14048}

\bibitem{xiao2024streamingllm}
Guangxuan Xiao, Yuandong Tian, Beidi Chen, Song Han, and Mike Lewis.
\newblock Efficient Streaming Language Models with Attention Sinks.
\newblock arXiv preprint arXiv:2309.17453, 2023. \url{https://arxiv.org/abs/2309.17453}

\bibitem{li2024snapkv}
Yuhong Li, Yingbing Huang, Bowen Yang, Bharat Venkitesh, Acyr Locatelli, Hanchen Ye, Tianle Cai, Patrick Lewis, et al.
\newblock SnapKV: LLM Knows What You are Looking for Before Generation.
\newblock arXiv preprint arXiv:2404.14469, 2024. \url{https://arxiv.org/abs/2404.14469}

\bibitem{cai2024pyramidkv}
Zefan Cai, Yichi Zhang, Bofei Gao, Yuliang Liu, Yucheng Li, Tianyu Liu, Keming Lu, Wayne Xiong, et al.
\newblock PyramidKV: Dynamic KV Cache Compression based on Pyramidal Information Funneling.
\newblock arXiv preprint arXiv:2406.02069, 2024. \url{https://arxiv.org/abs/2406.02069}

\bibitem{oren2024tova}
Matanel Oren, Michael Hassid, Nir Yarden, Yossi Adi, and Roy Schwartz.
\newblock Transformers are Multi-State RNNs.
\newblock arXiv preprint arXiv:2401.06104, 2024. \url{https://arxiv.org/abs/2401.06104}

\bibitem{chen2024sepllm}
Guoxuan Chen, Han Shi, Jiawei Li, Yihang Gao, Xiaozhe Ren, Yimeng Chen, Xin Jiang, Zhenguo Li, et al.
\newblock SepLLM: Accelerate Large Language Models by Compressing One Segment into One Separator.
\newblock arXiv preprint arXiv:2412.12094, 2024. \url{https://arxiv.org/abs/2412.12094}

\end{thebibliography}

\end{document}
